\subsection{Model descriptions}
The paper's Minecraft model has approximately 36 million parameters and a recurrent latent update built around a convolutional denoiser (Appendix D). The released \texttt{main} implementation uses temporal attention and samples with a sliding window of at most \texttt{n\_tokens}; this is distinct from the paper's continuous recurrent state. The README describes 100,000 training steps, whereas the released Minecraft checkpoint records \texttt{global\_step=110000}; the latter is the inspected artifact metadata. Independent noise levels let sampling schedules denoise tokens at different rates. Stabilization conditions future sampling on slightly noised past frames. The planning configuration uses a non-causal transformer of width 128 and 12 layers, with a pyramid schedule. The requested position-only checkpoints generate positions; a PD controller computes actions.
\subsection{Datasets}
The video input is the TECO Minecraft mini subset \cite{yan2023teco}, at $128\times128$ resolution with frame skip two and two context frames. The default validation configuration requests eight clips in one batch; eight clips each have been evaluated in the completed one-, two- and eight-times runs. Changing the horizon multiplier changes the dataset's clip enumeration, so comparisons across horizons are not guaranteed to use paired clips. Dataset padding has an explicit validity mask, which the instrumentation preserves. Maze evaluation targets D4RL \texttt{maze2d-medium-v1} and \texttt{maze2d-large-v1} \cite{fu2020d4rl} with 128 environments and configured episode lengths 600 and 800. Both maze environments pass reset/step checks, and both sparse-v1 datasets load through the repository. The registered environments use random starts and a fixed goal (\texttt{reset\_target=False}), rather than random goal pairs. The fixed upstream environment seed is retained.
\subsection{Hyperparameters}
The following values are verified in the completed one-, two- and eight-times video configurations; maze values are composed settings awaiting inference. Video uses $K=1000$, 100 DDIM sampling steps \cite{song2021ddim}, $\eta=0$, a sigmoid beta schedule, velocity prediction, frame stack one and stabilization level 15. Its configured training horizon is ten frames. Planned schedules are autoregressive with chunk size one and full-sequence with chunk size two, at multipliers $1,2,4,8$. Maze uses 50 DDIM steps, $\eta=0$, a linear schedule, clean-token prediction, frame stack ten, and open-loop horizon 50. Default guidance is three for medium and two for large; the proposed sweep is one through five. Resolved configurations are retained for each actual run, separately from planned configuration snapshots. Supplementary upstream FVD was disabled; this study measures consecutive-frame MSE and LPIPS with the specified analysis script.
\subsection{Experimental setup and code}
The upstream commit is \texttt{475e0bcab87545e48b24b39fb46a81fe59d80594}.

Local changes reside on \texttt{eecs553-draft}. The instrumentation records video tensors and validity masks, planner metrics and PNGs, and seeds Lightning before experiment construction. It does not alter model or sampler logic. The upstream planner fixes environment seeds to zero, so changing \texttt{DF\_SEED} varies model randomness without independently resampling environment seeds. Commands are generated by \texttt{scripts/run\_experiments.py}; offline runs use the dummy entity \texttt{offline-local} because upstream code rejects a null entity even offline. The command also sets \texttt{dataset.metadata} explicitly because its inherited interpolation references a missing \texttt{dataset.name}. Public fork: \textbf{[FORK URL]}.

Consecutive-frame MSE and AlexNet LPIPS \cite{zhang2018lpips} use pixels clamped to $[0,1]$, with unclamped ranges recorded separately. Per-frame CSV standard deviations describe variation across videos. Summary windows exclude context-only transitions; the first window ends before frame ten, while the last includes transitions ending in the last ten frames. Ground-truth transitions are included only if both validity entries are true. A nonzero, approximately flat curve near the reference distribution is consistent with absence of increasing temporal change, but cannot establish semantic correctness. A decline toward zero can indicate freezing; intensity and spatial standard deviation detect some saturation failures. Seed-level and video-level variation must be reported separately.
\subsection{Computational requirements}
The one-times inference took 440.02 seconds (0.12223 GPU-hours), or 55.00 seconds per video including startup and logging. The eight-times run took 6130.67 seconds (1.70297 GPU-hours), or 766.33 seconds per video. Sampled peak device memory was 15,827 and 14,343 MiB, respectively. The two-times run took 1243.28 seconds (0.34536 GPU-hours), with sampled peak 15,849 MiB. Completed inference totaled 2.17055 GPU-hours and metric postprocessing 0.00254 GPU-hours. Remaining runs are pending; these totals exclude live jobs. All released checkpoint SHA256 values are recorded in the artifact bundle. The host shell reports Python 3.14.9 and Conda is absent; a separate Python 3.10.22 virtual environment is now verified. The GPU is an RTX 4090 Laptop GPU with 16,376 MiB reported memory and driver 615.78.08. Initial device access failed under restricted execution; it succeeded after permissions changed, without driver modification. The run wrapper records timestamps, process wall time, CPU peak memory from GNU Time with \texttt{-v}, and sampled device-wide VRAM every two seconds. Device-wide sampled VRAM is not an exact process allocation peak. PyTorch 2.4.1 with CUDA 12.1, Torchvision 0.19.1 and Lightning 2.1.4 pass the runtime checks. Two synthetic GPU diagnostics took 0.001145 wall-GPU-hours in total, separate from research inference. A separate analysis environment supports document verification. GNU Time was built locally because the system binary is absent.
